\runde{3. Amtliche Erkenntnisse}

\mc{3.1 -- Der Knopf}
Im ET-Labor gilt die Taste \texttt{Autoscale} als sicherheitskritisch. Was passiert laut Gerüchten, wenn sie trotzdem gedrückt wird?
{Das Oszilloskop beantragt eine Weiterbildung in Digital Engineering.}
{Der Frequenzgang wird emotional und beginnt, den Raum zu bewerten.}
{Rolf Mettler kehrt aus dem Ruhestand zurück und stellt den Knopf auf manuell.}
{Die Zentralbahn fährt wegen EMV-Bedenken nur noch im Rechtecksignal.}
{C}
§
\offenfrage{3.2 -- Personalie}
Wer ist «the one and only»?
{Rolf Mettler.}

\mc{3.3 -- Protokollauszug}
In einer Besprechung fällt der Satz: «Man könnte das mitunter direkt
beantworten, oder zuerst überlegen, ob die Antwort nicht noch eine
andere Frage braucht. Oder» Wem wird die Aussage zugeordnet?
{Kilian Schuster}
{Thomas Hunziker}
{Christian Jost}
{Rolf Mettler}
{A}

\mc{3.4 -- Sprachstand}
Die Fachstelle für Mehrsprachigkeit untersucht eine Erklärung, die mit einem
Satz über Signalverarbeitung beginnt und als kurze, funktionierende Lösung
endet. Welche Sprache wurde darin am wahrscheinlichsten gesprochen?
{Python}
{Französisch}
{Latex}
{Die Sprache des Datenblatts}
{A}

\mc{3.10 -- Die Augenbrauenmethode}
Welche Reaktion bedeutet in einem analogen Labor mit hoher Wahrscheinlichkeit,
dass eine Erklärung noch einmal vollständig geprüft werden sollte?
{Eine sehr kleine Veränderung im Gesichtsausdruck der betreuenden Person.}
{Eine neue Version der Aufgabenstellung.}
{Ein besonders freundliches «passt schon».}
{Ein Screenshot des Messwerts.}
{A}

\runde{4. Fakultätsübergreifende Messungen}

\mc{4.1 -- Schwierigkeitsindikator}
Die ET-interne Kommission sucht einen verlässlichen Hinweis darauf, dass ein
Modul nicht sofort zur Flucht zwingt. Welcher Befund genügt?
{Auch WI-Studierende sitzen noch im Raum.}
{Die Aufgabenstellung ist auf Recyclingpapier gedruckt.}
{Der Beamer zeigt ein Balkendiagramm.}
{Jemand hat das Wort «Stakeholder» verwendet.}
{A}

\mc{4.2 -- Etymologie}
Im inoffiziellen Wörterbuch des wahren ETlers wird der Begriff
«Kanonenfutter» nicht erklärt, sondern mit einem einzelnen, fachlich
vollkommen ausreichenden Beispiel illustriert. Welches?
{Gebäudetechniker in ET2.}
{Die Beilage zum Mensamenü.}
{Ein Kollateralschaden beim Praktikumsaufbau.}
{Der Inhalt einer Schublade mit alten Bananensteckern.}
{A}
